\section{von Neumann double commutant theorem}

\begin{lemma}[Invariant subspaces and closure]
\label{lem:vN-closure-invariant}
\lean{VonNeumannDoubleCommutantTheorem.closure_invt}
\leanok
The topological closure of an invariant subspace of a continuous linear operator remains invariant.
\end{lemma}

\begin{lemma}[Orthogonal complements of invariant subspaces]
\label{lem:vN-orthogonal-invariant}
\lean{VonNeumannDoubleCommutantTheorem.orthogonal_invt}
\leanok
If a subspace is invariant under $a^*$, then its orthogonal complement is invariant under $a$.
\end{lemma}

\begin{definition}[Diagonal amplification]
\label{def:vN-diag}
\lean{VonNeumannDoubleCommutantTheorem.HSum, VonNeumannDoubleCommutantTheorem.diag}
\leanok
A bounded operator $a$ on $H$ acts diagonally on a finite Hilbert direct sum of copies of $H$.
\end{definition}

\begin{lemma}[Finite-vector approximation]
\label{lem:vN-finite-approx}
\lean{VonNeumannDoubleCommutantTheorem.finite_approx}
\leanok
\uses{lem:vN-closure-invariant, lem:vN-orthogonal-invariant, def:vN-diag}
If $x$ lies in the bicommutant of a star-subalgebra $S$, then on any finite tuple of vectors the diagonal action of $x$ lies in the norm closure of the corresponding diagonal actions of elements of $S$.
\end{lemma}

\begin{lemma}[Weak-operator approximation]
\label{lem:vN-wot-approx}
\lean{VonNeumannDoubleCommutantTheorem.wot_approx}
\leanok
\uses{lem:vN-finite-approx}
Every element of the bicommutant of $S$ lies in the weak-operator closure of $S$.
\end{lemma}

\begin{lemma}[Centralizers are weak-operator closed]
\label{lem:vN-centralizer-closed}
\lean{VonNeumannDoubleCommutantTheorem.wot_centralizer_closed}
\leanok
The image in the weak-operator topology of the centralizer of any set of bounded operators is closed.
\end{lemma}

\begin{theorem}[von Neumann double commutant theorem]
\label{thm:von-neumann-double-commutant}
\lean{VonNeumannDoubleCommutantTheorem.MainTheorem}
\leanok
\uses{lem:vN-wot-approx, lem:vN-centralizer-closed}
For a unital star-subalgebra $S\subseteq B(H)$, weak-operator closedness is equivalent to equality with its bicommutant.
\end{theorem}
